\section{Methodology: how the gap was established}
\label{sec:methodology}

A gap claim is only as good as the evidence behind it, and the easiest way to
get this wrong is to assert an absence from a cursory look at a portal. Three
independent methods were used, each implemented as a script that saves its
raw responses so that a reviewer can re-run it and diff the result.

\subsection{Method 1 --- specification audit (offline-reproducible)}

\code{scripts/probe\_gap.sh} audits Bhoonidhi's \emph{own published API
specification} rather than probing the live service. It parses the saved
specification document and mechanically enumerates:

\begin{itemize}
  \item which spatial parameters \ep{/data/search} documents;
  \item which spatial predicates are \emph{absent} --- \code{crosses},
        \code{touches}, \code{overlaps}, \code{disjoint},
        \code{s\_within} and \code{s\_contains};
  \item which CQL2 operations appear in the documented examples;
  \item whether any analysis or aggregation vocabulary appears anywhere;
  \item which endpoints exist, and whether any of them computes.
\end{itemize}

Auditing the specification rather than the service was a deliberate choice
with two benefits. First, it is the fairest possible reading: it tests what
NRSC says the API does, not what a particular deployment happened to answer
on a particular afternoon. Second, it is \textbf{reproducible offline} from
\code{docs/evidence/bhoonidhi-api-spec.html} and does not depend on NRSC being
reachable --- which mattered in practice, because both NRSC hosts stopped
answering from the build host partway through development (DNS resolved,
connections timed out, the rest of the internet was fine). The
\code{--refresh} flag re-fetches the specification when the hosts are back and
says so plainly when it cannot.

The script also deliberately searches for words that \emph{look} like
operators but are not. \code{contains}, \code{within} and \code{area} do
appear in the specification --- as English prose, never as parameters. Testing
for this explicitly turns a potential counter-example into part of the
argument instead of a hole in it.

\subsection{Method 2 --- comparative baseline (live)}

An audit alone establishes that the gap exists but not that it is
\emph{large}. On its own it invites the reply: \emph{server-side geospatial
computation is exotic, nobody deploys it, this is an unreasonable thing to
expect of a national portal.}

\code{scripts/probe\_baseline.sh} tests that reply against the field. It
probes ISRO's public stack for any processing endpoint, records how long the
relevant standards have existed, and then probes comparable Earth observation
portals live to see whether they actually expose the capability today. All
probes are single, unauthenticated, polite \code{GET}s against published
public endpoints. Portals that require registration are documented but
deliberately \emph{not} probed, on the grounds that hammering them from a
course project would be impolite rather than informative.

\subsection{Method 3 --- measuring the cost (live, against this service)}

The decisive method, because it converts an architectural complaint into a
number.

\code{scripts/measure\_redundancy.py} asks this project's own
\code{minimal-scene-set} process the question the catalogue cannot answer,
over an area-of-interest ladder from a taluk-sized window to a regional
transect, for both an optical and a radar collection. For each it records how
many scenes a search would return against how many are genuinely required to
cover the same ground to the same percentage.

\paragraph{What is measured and what is modelled.} This distinction is kept
explicit throughout, including in the script's own console output, because
conflating the two would be the easiest way to inflate the result:

\begin{itemize}
  \item the \textbf{scene counts are measured} live against the running
        service over real STAC footprint geometry. This is the part that
        carries the argument and it is not modelled;
  \item the \textbf{byte figures are modelled}: published nominal per-scene
        product sizes multiplied by those counts. The script never downloads
        imagery. A real order varies with processing level, tiling and
        compression;
  \item the ResourceSat figures are a \textbf{projection} onto Bhoonidhi's own
        holdings --- the measured footprint geometry is Sentinel-2 and
        Sentinel-1, so the per-scene volume is swapped for the NRSC product's
        while the geometry stays what was actually measured. Output labels it
        as such.
\end{itemize}

\paragraph{The measurement is a lower bound.} The footprint store holds 68
scenes over a three-month window, because that is what was harvested. Several
rows in the result are therefore capped by the store rather than by the area
of interest. A real Bhoonidhi search spans a multi-decade archive with many
revisits per year: it returns far more intersecting scenes for the same AOI
while the number genuinely needed to cover it stays roughly constant.
Redundancy therefore \emph{rises} with archive depth, and the figure reported
in Section~\ref{sec:cost} is the conservative end.

\subsection{What would falsify the claim}

Stated in advance, because a claim that cannot be falsified is not a finding:

\begin{itemize}
  \item if \ep{/data/search} documented any spatial operator beyond
        intersection, Method 1 reports \code{NOT MATCHED};
  \item if any ISRO endpoint answered a processing request, Method 2 reports
        \code{NOT MATCHED};
  \item if the minimal covering set were close in size to the candidate set,
        Method 3 would return a redundancy near zero and the cost argument
        would collapse.
\end{itemize}

Each script prints its own failure loudly rather than quietly passing. If
NRSC enables a processing endpoint, the honest response is to say the gap has
closed, and the scripts are written so that this is what happens.
